% Figura 1.1. Las dos tareas del diagnóstico dermatológico asistido.
% Se usa en sec:intro_contexto. El diagrama del pipeline completo (sec:met_pipeline) es la versión
% detallada de este mismo esquema: si uno de los dos cambia, hay que revisar el otro.

\begin{figure}[htbp]
    \centering
    \begin{tikzpicture}[x=1mm, y=1mm]

        \node[figentrada, text width=27mm] (img) at (12, 0) {Imagen\\dermatoscópica};

        \node[figtarea, text width=36mm] (t1) at (58, 14)
            {\textbf{Clasificar}\\asignar una categoría};
        \node[figtarea, text width=36mm] (t2) at (58, -14)
            {\textbf{Describir}\\redactar el texto\\de la lesión};

        \node[figsalida, text width=34mm] (cat) at (114, 14)
            {Categoría diagnóstica\\[2pt]{\footnotesize de ella se deriva la\\detección de melanoma}};
        \node[figsalida, text width=34mm] (txt) at (114, -14)
            {Reporte clínico\\[2pt]{\footnotesize el texto que\\describe la lesión}};

        \draw[figflecha] (img.east) -- ++(6, 0) |- (t1.west);
        \draw[figflecha] (img.east) -- ++(6, 0) |- (t2.west);
        \draw[figflecha] (t1.east) -- (cat.west);
        \draw[figflecha] (t2.east) -- (txt.west);

        \begin{scope}[on background layer]
            \node[figgrupo={lo que un modelo visión--lenguaje promete resolver con un solo modelo},
                  fit=(t1) (t2)] {};
        \end{scope}

    \end{tikzpicture}
    \caption[Las dos tareas del diagnóstico dermatológico asistido]{Las dos tareas del diagnóstico
    dermatológico asistido. De una misma imagen se espera una categoría diagnóstica y un texto que describa
    la lesión, y la detección de melanoma es la tarea binaria que se deriva de esa clasificación. Un reporte
    es correcto solo si se cumplen las dos condiciones a la vez: que la categoría sea la acertada y que la
    descripción corresponda a la imagen. Las dos condiciones son independientes y de ahí que se evalúen por
    separado.}
    \label{fig:intro_dos_tareas}
\end{figure}
